\section{No selective effect found}
\label{sec:causal-results}

\subsection{Where the tool choice is carried}
\label{sec:map}
Before steering, we mapped where the choice of \texttt{answer} is made. On two \fsp{} and two \fsm{} completion steps, we patched the residual stream from the same task's previous step (whose recorded action is a click) at each of \cfgBigLayers{} layers and four positions, \cmProbes{} patches in all, with the original reasoning held fixed. The quantity is the logit margin between the first token of \texttt{answer} and that of the best other tool. Identity patches give exactly zero.

Because the reasoning is forced to its recorded text, a patch before or inside it cannot change what the model writes, so the decision is invisible at the end of the prefill and barely visible mid-reasoning by construction: the largest mean margin change is \cmPrefillPos{} and \cmMidPos{} logits there, for \fsp. It appears only at the tool-name position, from about layer 16 upward, where patching moves the margin by up to \cmMinPos{} logits on \fsp{} and \cmMinNeg{} on \fsm{} (Fig.~\ref{fig:map}). Component scans single out three attention heads, L19/H15, L23/H1 and L31/H14, and ten late MLP neurons. Their attention falls mostly on the system-prompt definition of \texttt{answer} (L19/H15, L23/H1) and on the JSON being written (L31/H14). % source: data/pod_astra/circuit/REPORT.md "Où regardent les têtes ?"
Patching all of them together from the donor step lowers the margin by \cjDeltaMin--\cjDeltaMax{} logits, and flips \cjFlips{} of the \cjN{} decisions (smallest remaining margin \cjMarginMin). Across all \cmCompProbes{} component patches and \cjAll{} joint patches, \cmCompFlips{} and \cjAllFlips{} decisions flip: the components move the margin, not the action. The only flips come from residual patches (\cmFlipsPos{} on \fsp{} and \cmFlipsNeg{} on \fsm{} of \cmProbes), at the tool name in the last layers, where replacing the residual stream transplants the donor's choice into the output head. They show where the choice is carried, not where a false judgement is formed. % source: data/pod_astra/circuit/REPORT.md "Carte causale"
 The effect is the same on true and false completions (Fig.~\ref{fig:map}c). On these four contexts, one of them infeasible, and with component scans limited to layers 19, 23, 30 and 31, these components behave like a tool-name retrieval circuit, not a false-success circuit. % source: data/pod_astra/circuit/REPORT.md (fine scans on L19/23/30/31)

\paragraph{A ceiling test.} The patches above replace one layer at a time; a ceiling test asks how much must be replaced at once. With the recorded input forced and only logits read, the outputs of whole blocks of writers are swapped in at the decision position from a donor step where the model clicked: the same task's previous step for two \fsp{} contexts, another task's step for one true completion (Fig.~\ref{fig:ceiling}). On the first \fsp{} context (margin \ceilBase), all 32 MLPs together bring the margin to \ceilMlpAll{} without flipping it, while the MLPs and token mixers (attention and DeltaNet) of layers 12--23 flip it (\ceilBothMiddle), as do those of layers 24--31 (\ceilBothLate), and those of layers 0--11 do not (\ceilBothEarly). The second \fsp{} context behaves the same (MLPs \ceilSecMlpAll, block 12--23 \ceilSecBothMiddle), and so does the true completion (block 12--23: from \ceilTrueBase{} to \ceilTrueBothMiddle). Swapping with itself and swapping the final residual reproduce the baseline and the donor exactly. % source: data/pod_astra/broad/ceiling.jsonl, ceiling_resume.jsonl, ceiling_checks.json (asserted in figs/broad.py)
Three caveats bound this result. The first \fsp{} context is an infeasible task, on which \texttt{answer} (with a refusal) is the correct tool (Appendix~\ref{app:infeasible}). The donor differs in screen, reasoning, thought and position, not only in decision. And since a true completion flips too, the test locates where the identity of the next tool token is assembled, not a mechanism of false success. No particular block is needed: the MLPs and mixers of layers 12--23 or of layers 24--31 each flip the first context, and removing a single direction of the swap difference from the writers of layers 12--23 brings the margins of the two \fsp{} contexts to \projOneFspA{} and \projOneFspB{}, but also those of two true completions to \projOneTrueA{} and \projOneTrueB. % source: data/pod_astra/broad/projection_summary.json (rank 1, dose 1)
 It is consistent with the failure of every edit confined to one or two layers or to the MLPs, and it says nothing about which edit would separate false from true completions.

\subsection{The tool call leans on the text already written}
\label{sec:anchor}
The tool call leans on the text the model has already written, for true completions as for false ones. On the \paFsN{} false and \paTrueN{} true completions of the strictly feasible cohort (\S\ref{sec:levers}), teacher-forced with the screen, history, reasoning and token count fixed, replacing the note and thought by a neutral text of the same token length (``Observing the current screen.'') removes \paNeutralFs{} logits of \texttt{answer} margin on false completions (95\% CI \paNeutralFsCi; original \paOrigFs) and \paNeutralTrue{} on true ones (\paNeutralTrueCi; original \paOrigTrue), \paFracFs\% and \paFracTrue\% of the margin per task. The two losses differ by \paInter{} logits (\paInterCi; label permutation $p=\paInterP$). A counterfactual text of the same length (``The task is not complete yet; I still need to inspect the screen.'') leaves \paCounterLeftFs{} and \paCounterLeftTrue{} logits, nearly cancelling the margin on average (\texttt{answer} still leads in \paCounterPosFs/\paFsN{} and \paCounterPosTrue/\paTrueN{} contexts), while neutralising the recorded reasoning instead removes only \paReasonFs{} and \paReasonTrue. % source: data/pod_astra/proofanchor/RESULT.json anchoring, PROTOCOL.md conditions a-d, VERIFY_FINAL.md (figs/proofanchor.py)
One true completion has an empty note; without it, the true-completion loss is \paNeutralTrueSens. The written note and thought thus carry \paFracTrue--\paFracFs\% of the \texttt{answer} margin in both strata, yet without them \texttt{answer} remains the preferred tool in \paNeutralPosAll{} of \paN{} contexts (\paNeutralPosFs/\paFsN{} false, \paNeutralPosTrue/\paTrueN{} true): the conclusion is reinforced by the note, not created by it, although the success claim in the retained reasoning may carry the rest. This is an attenuation of the margin, not a change of decision. Nor does a probe find a robust discriminator in the residual stream of these contexts: of \paNProbes{} out-of-fold probes (\paNLayers{} layers, three positions, four models), the best, chosen after the fact, reaches AUROC \paBestAuc{} (\paBestCi) against \paStepAuc{} for the raw step index, and its permutation $p$ of \paPermP{} becomes \paPermPadj{} after Bonferroni correction for the \paNProbes{} probes. With 1{,}000 permutations the corrected $p$ cannot fall below 0.072, so this corrected test had no power at all. With these sample sizes a pre-specified score would need an AUROC of at least \paPowOne{} (one-sided) to \paPowTwo{} (two-sided) to be detected with 80\% power, and about 0.85 after a correction for \paNProbes{} probes, % source: data/pod_astra/proofanchor/RESULT.json permutation.minimum_possible_adjusted_p 0.0719, power.actual_1000_permutation_bonferroni_power_at_05 0.0, power.minimum_auroc_for_80_percent_power.familywise_72 0.845
 so the absence of the information is not proven. These are artificial inputs, and the joint replacement does not separate the wording of the conclusion from the details the note describes.

A first version of the test, on two \fsp{} contexts, illustrates it with a donor's text rather than a neutral one: it keeps the screen, the history and the weights, and replaces only the text the model has generated at the decision step with that of the same task's previous step, where the model clicked. On the first \fsp{} context (margin \ancAOrig), the donor's reasoning alone lowers the margin to \ancAReason, the donor's JSON prefix alone (note and thought) to \ancAJson, and both together flip it (\ancABoth). The second context behaves the same (\ancBOrig; reasoning \ancBReason, JSON prefix \ancBJson, both \ancBBoth). % source: data/pod_astra/broad/anchoring.jsonl, anchoring_summary.json, anchoring_plan.json
Three limits weighed on that version, the first of which the cohort test above removes: it had no control on true completions or with a neutral text. Only its second context is feasible, and the donor is the previous step, whose text announced a click, so that copying its text largely copies its decision. And the text does not account for everything: with all of it swapped, the margin stays \ancAScreen{} and \ancBScreen{} logits above the donor's, a share that comes from the screen and the history. Within these limits, the tool call depends causally on the text already written, the note and thought weighing most, and the \margPosMin--\margPosMax-logit margin of \S\ref{par:ceiling} is largely conditioned on that text rather than an intrinsic robustness of the decision.
An attribution on the same cohort locates the channel. Replacing the inner tokens of the note and thought with spaces, at the same length and positions and with the screen, history and reasoning kept, lowers the teacher-forced margin by \mlNeutralFs{} logits on false completions (mean \mlBaseFs{} before) and by \mlNeutralTrue{} on true ones (mean \mlBaseTrue{} before, measured on its own cached replay path), less than the neutral text above, and lowers it on every context of both (\mlNeutralFsPos/\mlFs{} and \mlNeutralTruePos/\mlTrue); the difference between the two losses is \mlNeutralSel{} logits (95\% CI \mlNeutralSelCi). The note and thought carry the decision as much for true as for false completions: the channel is real, but not selective. Two kinds of patch probe the route (Fig.~\ref{fig:moonleak}). Replacing the output of a single DeltaNet layer at the decision position by its value under the neutralised text costs \mlDeltaTopMin--\mlDeltaTopMax{} logits at layers \mlDeltaTopLayers{} (zero-based) and at most \mlDeltaRestMax{} at the other layers. Blocking the final query's attention to these tokens, one head at a time over the \mlHeads{} heads of the eight full-attention layers, costs at most \mlHeadMaxFs{} logits on false completions and \mlHeadMaxAny{} on true ones, and at best \mlBestHead{} logits more on false than on true. The two tests are not symmetric: a whole-layer output patch includes every upstream effect on that layer, while a head block cuts one attention link of the final query, and no whole full-attention layer was patched as a control. So single DeltaNet layer outputs carry \mlDeltaTopMin--\mlDeltaTopMax{} logits each and blocking the final query's attention one head at a time removes at most \mlHeadMaxAny; this does not show that the recurrent memory rather than attention carries the conclusion. The same contexts define the comparison and the cohort, so this screen is exploratory, and layer-output patches locate carriers of the change, not its mechanism. % source: data/pod_astra/moonleak/attribution_summary.json delta, heads; independent_attribution_stats.json (figs/moonleak.py)

\placefig{fig-map}

\placefig{fig-ceiling}

\placefig{fig-moonleak}

\subsection{Every steering arm}\label{sec:every-arm}
Table~\ref{tab:causal} lists every intervention run on the 4B false-success target, each under the gate of \S\ref{sec:protocol} and on six screening contexts, the same six for every arm except the neuron and \texttt{perceive} arms, which use their own. Each line tests a distinct hypothesis; four come from the literature, the others from our own analyses.
\begin{itemize}
\item \emph{Tool choice is linearly steerable} \citep{wu2026toolcalling}, tested with the position-matched \fsp{}$-$\fsm{} direction rather than a tool-choice direction: add it at the last prefill token with coefficient $\alpha=\cAlphaSixteen$ at layer 16 and $\alpha=\cAlphaTwentyFour$ at layer 24 (absolute magnitudes; a negative $\alpha$ subtracts the direction, and Table~\ref{tab:causal} gives the dose relative to the residual norm).
\item \emph{The decision is taken during reasoning} (\S\ref{sec:think}): subtract a ``conclusion'' direction fitted on the reasoning tokens ($\alpha=-0.5$) at layer 31, at every generated token, including the JSON after \texttt{</think>}, so the arm can also act on the tool choice directly. % source: data/pod_astra/think/REPORT.md "Portée mécanistique"
\item \emph{False success is excess expected value}, a hypothesis derived from \citet{jiang2026value} and \citet{kumaran2026commit}: push toward low value or low commitment at layer 24 or 18.
\item \emph{Late MLP components write the decision} \citep{dutta2026termination}, tested here at the neuron level: zero the top-16 or top-64 neurons by direct attribution to the \texttt{answer} logit.
\item \emph{The model fails to check}: add a ``verify'' versus ``conclude'' direction (contrastive activation addition, \citealp{rimsky2023caa}) at layer \verifyLayer{} at 1 and 2 times the residual norm, over the prefill and the first 20 generated tokens.
\item \emph{False success is a perception of completion}: subtract the imminent-completion direction of \S\ref{sec:perceive} at layer \pcLayer.
\item \emph{The model opens with its conclusion} (\S\ref{sec:think}): raise the output row of the token ``Wait'' for the states that open a completion, a \waitWeights-weight edit of the output head.
\item \emph{Memory of the early episode fades}: lengthen or shorten the retention of the nine late linear-attention layers by a factor 1.25.
\item \emph{The circuit of \S\ref{sec:map} carries it}: scale its three heads and ten neurons (C1) by $1-\alpha$ with $\alpha=0.5$ and $1$, against a random direction of the same contribution norm.
\end{itemize}
No arm passes the gate, and none shows a large effect. Most leave the false-``done'' rate at its original value, including removing the whole C1 circuit ($\alpha=1$), which answers on every false completion, like the random control. The neuron arms, screened on six other disjoint contexts whose own baseline is \cNeuronBase{} \fsp{} answers, were stopped by a mathematical bound: with \cNeuronObs{} answers already observed on the first unseen \fsp{} task, at most \cNeuronMax{} answers out of 8 would have been admissible. The value-axis arm did change the text. An LLM-agent reader found an explicit intention to re-evaluate in \doubtPushFsPos{} \fsp{} reasoning traces against \doubtBaseFsPos{} at baseline, but \doubtPushStrict{} verification actions and the same answer rate. % source: data/pod_astra/value/*_VERIFICATION_ANNOTATIONS.json (secondary/strict flags)
The model wrote that it should check, then answered anyway.

The ``Wait'' row is the one edit that lowers false answers without touching anything else: \waitFsp{} on false completions, with \waitFsm{} true completions kept, \waitHealthy{} healthy actions valid and \waitText{} general-text controls token-identical. It still fails the 50\% relative criterion, and an LLM-agent review rates none of the two replacement actions useful: one reloads a settings page for an invented reason, the other scrolls without fixing the dates. % source: data/pod_astra/mit1/REPORT.md "Lot6 wait2"; data/pod_astra/mit1/semantic_lot06.json
Making the model pause is not enough to make it repair.
A conditional variant adds an early end-of-reasoning row to the same edit (\earlyExitWeights{} weights). Inside its construction set, \hfourInFalse{} false completions still answer. On one Budget sample the model reasons correctly that next month is December, says the dates must be changed, and generates a click on the date summary of the results page; it was not executed. % source: data/pod_astra/mit1/semantic_lot14.json; data/pod_astra/mit1/REPORT.md first section
These are intentions to correct, not corrections: an LLM-agent review of the screenshot finds that the dates sit in plain summary text, next to a separate ``Modify Rental Details'' control, and that no step proves the summary opens an editor. All \hfourIntentions{} such generated clicks, over \hfourIntentLots{} lots of this edit and its diagnostics, are reclassified as unverified, \hfourConfirmed{} confirmed useful; no effect on the interface or on the OSWorld score is shown. % source: data/pod_astra/mit1/review_budget_affordance.{json,md} (sha 0aa1bc6a); data/local/JOURNAL.md entries 00:15 and 00:20 UTC
Nor does the edit transfer. On nine training tasks left out of its construction (8 of the 12 \fsp{} outputs come from two infeasible tasks), the final audit counts \hfourOutUseful{} useful actions on false completions (original \hfourBaseUseful), \hfourOutErrId{} outputs that name the real error (original \hfourBaseErrId), and \hfourOutFalse{} false finishes (original \hfourBaseFalse), while true completions go from \hfourBaseTrue{} to \hfourOutTrue{}. % source: data/pod_astra/mit1/semantic_lot19.json (sha 44a9011b; fs+ tasks gimp/8ea73f6f and os/fe41f596 infeasible by evaluator func, 4 outputs each; supersedes the first audit's 1/12 useful and 5/12 identified); SEMANTIC_AUDIT_LOT15_ADDED_BASELINE.json
The edit leaves a mechanistic signal, the model recognising the error then rationalising it away, and no demonstrated useful effect.

\paragraph{These screens are pilots.} Each efficacy cell holds 8 generations from two tasks. 8/8 has a Wilson 95\% interval \citep{wilson1927} of [\wilsonEightLo, 100]\%, and samples clustered in two tasks make the real uncertainty larger, so a rejection means only that no large effect appears on these contexts, not that none exists. By design the gate flags a reduction to at most \gateMaxAnswers/8 answers; smaller effects pass unnoticed. Table~\ref{tab:causal} gives the interval of every efficacy cell.

\paragraph{A ceiling effect is likely.}\label{par:ceiling} At the tool-name position the margin of \texttt{answer} over the best other tool is \margPosMin--\margPosMax{} logits on false completions and \margNegMin--\margNegMax{} on true ones. This position comes after the JSON \texttt{thought} field, which the output schema writes before the tool call and which already states that the task is complete, so the margin measures how confidently the model carries forward a decision it has already verbalised, not the margin of the decision itself. % source: ~/osw/rtg/Agent/holo_agent/tools.py make_step_model (note, thought, tool_call); data/local/reports/REDTEAM_0305.md point 3
Even so, the perturbations tested are small: next to the residual stream, they are \doseSixteen{} (layer 16) and \doseTwentyFour{} (layer 24) of the mean residual norm for the matched direction, \doseThink{} for the reasoning-token direction (difference-of-means norm \thinkScale{} against a residual norm of \thinkRThirtyOne{} at layer 31), and 0.5 for the value axes, % source: data/pod_astra/value/REPORT.md "Norme = norme résiduelle moyenne"
while patching the whole C1 circuit from the donor step lowers the margin by \cjDeltaMin--\cjDeltaMax{} logits without flipping it. Larger doses do not break the ceiling either. The ``verify'' direction at one and two times the residual norm leaves \fsp{} at 8/8; at dose 2 a healthy context loops on ``I see'' up to the 3{,}072-token cap, which stopped the arm after \verifyRejCompleted{} of 24 outputs. % source: data/pod_astra/verify/verify_a2.rejected.json; JOURNAL.md 22:10 UTC astra-verify
Dose sweeps at layers 24 and 31 with random controls were not run.

A second target, getting stuck at the step limit, shows the opposite failure of steering. A residual addition at layer 16 or 24 damaged healthy actions before any stuck context was generated; after \stuckSixteenObs{} outputs, the best reachable health rate was \stuckSixteenBest\%, below the required \stuckSixteenReq\%. % source: data/pod_astra/stuck/CAMPAIGN_STATUS.json

\placefig{tab-causal}

\subsection{A detector of imminent completion beats position, not false success}
\label{sec:perceive}
All directions above contrast false and true completions, which differ in position. A separate analysis sidesteps this by learning a direction \emph{inside successful episodes}: the last-prefill residual of the final, correct \texttt{answer} step minus that of the step just before it, on \pcPairs{} pairs from \pcTasks{} construction tasks (reserve excluded, cross-validation grouped by task across both passes). Because both members of a pair sit next to each other in the same episode, position cannot separate them across tasks. The direction peaks at layer \pcLayer{} with task-grouped AUROC \pcAuc, against \pcBaseStep{} for the step index and \pcBaseCtx{} for context length (Fig.~\ref{fig:perceive}a). This is the first signal in this study that position does not explain, in contrast to the false-success direction of Table~\ref{tab:fs-auc}. Within a pair, both baselines are trivially perfect, so the pooled AUROC is the relevant comparison, and the layer was chosen on the same cross-validation.

Read on \pcMatchN{} false/true completion pairs matched on pass, step and context length, with each task's direction fitted without that task (3 of the 14 false completions are infeasible tasks, and one false-completion task appears in two pairs), the direction rises at the ``done'' step in \pcFsPosPositive/\pcMatchN{} false and \pcFsNegPositive/\pcMatchN{} true completions (mean \pcFsPosMargin{} and \pcFsNegMargin{} standard deviations; Fig.~\ref{fig:perceive}b). At the terminal step, false completions sit \pcTermDiff{} standard deviations from true ones, above them in only \pcFsPosGt/\pcMatchN{} pairs. It is therefore a detector of imminent completion, not of false success: it fires before a false ``done'' as it does before a true one. On these \pcMatchN{} pairs, the projection reads a false completion as complete. This is compatible with a perception error: the screen is taken as showing a finished task. Two same-task image pairs from construction data illustrate it, a spreadsheet cell claimed green whose text has no green pixels, and rental dates claimed for next month while the page shows the current one. % source: data/pod_astra/perceive/REPORT.md "Screen and visual evidence"
It does not separate misperception from commitment to the action. Subtracting it at layer \pcLayer{} (up to \perceiveMaxDose{} of the residual norm, on its own screening contexts whose reference is \perceiveFsBase/8 \fsp{} answers; both \fsp{} contexts are infeasible tasks, so this screen tests no feasible false success) gives \perceiveFsSeries{} \fsp{} answers out of 8 for $\alpha=0.5,1,2,4$, against \perceiveFsRandom/8 for a random direction of the same norm, and no dose passes the gate (Table~\ref{tab:causal}). The few drops are not repairs: the non-answer actions that replace \texttt{answer} are empty or invented text entries, such as a made-up e-mail address. % source: data/pod_astra/perceive/REPORT.md "Génération causale"
Attention maps on two same-task pairs found no head, out of \perceiveHeadsTested{} tested, that reads the evidence region more when the model succeeds (\perceiveHeadsEligible/\perceiveHeadsTested{} pass the frozen rule). The perception hypothesis is not refuted by this small screen, but it has no support beyond the projection of Fig.~\ref{fig:perceive}.

\placefig{fig-perceive}

\subsection{Weight edits: levers without established selectivity}
\label{sec:levers}
Weight edits that need no direction, or that target a mechanism, are detailed in Appendix~\ref{app:edits}; we summarise them here. Extrapolating the fine-tune away from Qwen, $W(\lambda)=W_{\text{Qwen}}+\lambda\,(W_{\text{Holo}}-W_{\text{Qwen}})$, called ExPO \citep{zheng2024expo}, leaves false answers at 16/16 in every arm but one (\expoFspMin/16, MLPs only, where true completions drop more); the edit at $\lambda=1.25$ keeps text and JSON intact and runs online (\S\ref{sec:online}). The mechanism-level edits H2--H10 (memory gates, MLPs or all language tensors moved toward Qwen, the end-of-thinking embedding, the attention temperature, the rotary share of the query gain) leave the tool call in place, break generation, or move the call only on the infeasible context, where theme clicks replace an answer that should have explained the impossibility (Table~\ref{tab:attractor}); two-neuron circuits, a narrow MLP edit adapted from weight-space programming \citep{takehi2026wsp} and a corrective instruction injected into the reasoning do no better. Two edits do move the decision on feasible contexts, and neither shows an established selectivity.

Merging in a donor trained on desktop tasks moves the decision globally, not selectively. Adding a fraction $\beta$ of the task vector \citep{ilharco2023task} of CUA-Lite, a Qwen3.5-4B fine-tune on ScaleCUA desktop data, to the language layers of Holo ($W_{\text{Holo}}+\beta\,(W_{\text{CUA}}-W_{\text{Qwen}})$; vision tower unchanged; input embeddings merged, the output head being tied to them at load time), and serving it with the same 51 calls as the exact original model, re-served on the same GPU type, image and vLLM version (on a separate Modal account), lowers false-completion answers to \mergeCuaQFsp{} at $\beta=0.25$ and \mergeCuaFsp{} at $\beta=0.5$, and true-completion answers further, to \mergeCuaQFsm{} and \mergeCuaFsm{} (original \mergeCuaOrigFsp{} and \mergeCuaOrigFsm), with every healthy action valid and no invalid output; on the feasible \fsp{} context alone, $\beta=0.5$ answers \infMergeCuaHalfExcl. The dose--response is a global drop of \texttt{answer}, not a discrimination. An LLM-agent audit finds the replacing actions coherent and unnecessary: scrolls after the model has recognised success, and verification clicks after true completions. The model continues instead of concluding, whatever the truth of the completion. A donor trained on filtered web tasks, Fara1.5-4B, leaves the decision in place at $\beta=0.5$ (\mergeFaraFsp{} and \mergeFaraFsm). BF16 rounding blurs small merges: at $\beta=0.5$, \mergeCuaUnchanged{} of the \mergeCuaLangTensors{} language tensors are left unchanged and the update norm is \mergeCuaActualNorm{} against an intended \mergeCuaIntendedNorm; at $\beta=0.25$, \mergeCuaQUnchanged{} tensors are unchanged. CUA-Lite's training data may overlap OSWorld tasks. % source: data/pod_astra/merge/cua_metrics.json per_tensor (embed_tokens merged), data/pod_hf/config_holo31_4B.json tie_word_embeddings, REPORT_orbit.md (lm_head kept in the file); data/pod_astra/merge/STATIC_RESULTS.json; VERIFY_MODAL_REFERENCE.{md,json} (original revision, H100, vLLM 0.23); VERIFY_CUA_B050_BEHAVIOR.{json,md}, VERIFY_CUA_B025_BEHAVIOR.md, VERIFY_FARA_B050_BEHAVIOR.md; REPORT.md (donor origin, contamination limitation)
Online, the merge shows no detectable gain over an original served on the same GPU type, image and vLLM version, from a separate Modal account: \onAllCuaDelta{} points over \onAllCuaN{} tasks (95\% CI \onAllCuaCi; \S\ref{sec:online}).

A closed-form edit of the writers that the ceiling test points to moves the decision as well, unselectively on the screening contexts; on a larger feasible cohort its selectivity is not established either way (below). In the spirit of AlphaEdit \citep{fang2024alphaedit}, the \alphaNWriters{} residual writers of layers \alphaLayers{} (MLP down-projections and mixer output projections) receive a ridge update that maps \fsp{} keys (constructed on the ceiling-test contexts, the first of them infeasible) toward the donor displacement while being projected off the span of protected keys (true completions, healthy steps and text prompts). On contexts outside its construction, served next to the original on the same GPU, false-completion answers fall from \alphaOrigFsp{} to \alphaEditFsp{} but true-completion answers fall further, from \alphaOrigFsm{} to \alphaEditFsm, with \alphaEditHealthy{} healthy actions valid and \alphaEditInv{} invalid outputs. The whole \fsp{} drop comes from one of the two contexts (\alphaOrigFspCtxB{} to \alphaEditFspCtxB; the other stays at \alphaEditFspCtxA). A first visual audit of the replacing actions, which had described the screenshot of the previous step, was retracted; the corrected audit finds that on the e-mail true completion all four recipients were already complete, so the three generated \texttt{write} continuations there would potentially damage a correct state (they were not executed). % source: data/pod_astra/broad/{eval_original,eval_p100}/summary.json, derive_residual.json; data/local/JOURNAL.md 02:14, 02:35 and 02:37 UTC (semantic_p100.{json,md} is the retracted audit)
A random update on the same writers, with the same right factors and its per-writer BF16 perturbation norm matched to the learned one within 0.1\%, leaves the screening contexts in place (\fsp{} \alphaRandFsp, true completions \alphaRandFsm, healthy \alphaRandHealthy): on these contexts the learned edit moves stopping and this one random draw does not, without selectivity. % source: data/pod_astra/broad/eval_random100/summary.json; data/pod_astra/broad/REPORT.md (random controls)
At half the dose the edit gives \alphaHalfFsp{} and \alphaHalfFsm: at both doses, true completions fall on these contexts as fast as or faster than false ones. % source: data/pod_astra/broad/eval_p050/summary.json
A 25\% dose was not run.

Two distinct methods, a merge and a closed-form edit, thus move the decision. On the screening contexts, true completions fall at least as much as false ones in both, on two true-completion contexts; on the feasible cohort below, this holds for the merge, while the selectivity of the closed-form edit is not established. A simple argument makes a lever without a discriminator unsurprising. An edit applied at or after layer $\ell$ acts on a function of the residual stream at layer $\ell$ (all positions, and the recurrent state of the linear-attention layers), so it can treat a false and a true completion differently only to the extent that some function of that stream separates them, and to be useful that function must generalise to tasks it was not built on. Our probes, which read the stream at two positions, found no such separation beyond episode length (Table~\ref{tab:fs-auc}).
A direct test on the writers that both methods edit finds the same. On \rmFsp{} false and \rmFsm{} true completions outside the reserve and the evaluation contexts, a ridge probe on the inputs (keys) of the \rmWriters{} writers of layers 12--23, cross-validated by task, does no better than position. Averaged within test folds and weighted by their \rmPairs{} positive--negative pairs, its median AUROC is \rmPreAuc{} (range \rmPreRange) at the end of the prefill and \rmEndAuc{} at the end of reasoning, against \rmBaseStep{} for the same probe on the step index; pooled out of fold, it is \rmPreOof{} and \rmEndOof{} against \rmBaseStepOof. % source: data/pod_astra/regmean/overlap.json (pair_weighted_fold_auc and oof_auc, cv_probe in astra-regmean/analyze.py), overlap_summary.json; manifest.json (sha f517da46)
This pool is weaker than it should be: \rmFspInfeasible{} of its \rmFsp{} false completions are infeasible tasks, on which \texttt{answer} is the correct tool, so the probe mixes in cases where the label marks no wrong decision, which weakens the argument. The strictly feasible cohort below avoids this for behaviour, though not for the probe. % source: data/pod_astra/power/BUILD_AUDIT.json, PROTOCOL.md; data/local/osworld_tasks/*.json evaluator.func
The two key subspaces are not identical (median principal angle \rmPreAngle--\rmEndAngle$^\circ$), but \rmEndEnergy--\rmPreEnergy\% of the false-completion key energy lies in the true-completion subspace. This is an empirical, non-total overlap together with a strong confound of position and length, not an impossibility: a low AUROC for a linear probe is not an information bound, the probes do not exhaust the stream, and even identical subspaces would not force a selective edit or merge to act globally, which would require proportional covariances. % source: data/pod_astra/regmean/overlap_summary.json, overlap.json (its own caveat)

The agent's own text carries no more. On \tsNFs{} feasible false and \tsNTrue{} true successes from \tsNTasks{} tasks, the best of \tsNModels{} detectors reading the final note, thought and reasoning, chosen after the fact, reaches a task-grouped AUROC of \tsBestAuc{} (95\% CI \tsBestCi) against \tsStepAuc{} (\tsStepCi) for the final step index; within the \tsMixed{} tasks that end both ways, every detector lies between \tsMixedMin{} and \tsMixedMax, and a within-task permutation test gives $p=\tsPermP$. Only the hesitation between samples carries a little: over four samples per context of the feasible cohort below, the original answers at a mean rate of \tsConsRateFs{} on false and \tsConsRateTrue{} on true completions, an AUROC of \tsConsAuc{} (\tsConsCi). Linear probes at two positions found no separation beyond length, and the text detectors none beyond the step index; only the hesitation from one sample to the next carries a little; at these sample sizes this is an absence of evidence, not a proof of absence. % source: data/local/text_signal/RESULT.json, REPORT.md (figs/closures.py)

\paragraph{A strictly feasible cohort.} To take infeasible contexts out of the evidence, the original, the merge, the closed-form edit and its random control were re-measured on a cohort fixed before any output: \pwOrigTasks{} feasible false completions, \pwFsm{} true completions and \pwHealthy{} healthy contexts, all from distinct tasks whose evaluators were read directly, with $k=4$ paired seeds per context and an LLM-agent audit of every count. % source: data/pod_astra/power/RESULT.json, PROTOCOL.md, BUILD_AUDIT.json, VERIFY_FINAL.json, REPORT.md (figs/power.py)
The original answers on \pwOrigFsp{} false and \pwOrigFsm{} true completions. The merge answers on \pwCuaFsp{} and \pwCuaFsm: false answers change by \pwCuaDeltaFsp{} points (95\% CI \pwCuaCiFsp), true ones by \pwCuaDeltaFsm{} (\pwCuaCiFsm), a selectivity (first difference minus second) of \pwCuaSel{} (\pwCuaSelCi). On feasible tasks, the merge is a lever without a discriminator. The closed-form edit answers on \pwBroadFsp{} and \pwBroadFsm: against the original, false answers change by \pwBroadDeltaFsp{} points (\pwBroadCiFsp), on \pwBroadFspDown{} of the \pwOrigTasks{} tasks rather than one context, and true ones by \pwBroadDeltaFsm{} (\pwBroadCiFsm), a selectivity of \pwBroadSel{} (\pwBroadSelCi). Its point estimate favours false completions, but the interval includes zero: neither a selective repair nor its absence is established. The drop survives removing the three tasks that overlap its construction or preservation sets (\pwBroadSensDeltaFsp{} points on \pwSensFsp{} false completions, \pwBroadSensCiFsp). Its random control, a single draw that left the screening contexts in place, here raises false answers by \pwRandDeltaFsp{} points (\pwRandCiFsp) and leaves true ones unchanged (\pwRandDeltaFsm); the learned edit lies \pwBroadVsRandDeltaFsp{} points below it on false completions (\pwBroadVsRandCiFsp), but because the control itself moves the result, the contrast with the original above is the relevant one. Two other perturbations that do not target the decision also raise false answers on this cohort. Setting to zero the four attention heads ranked most selective in the attribution of \S\ref{sec:anchor} changes them by \mlZeroFsp{} points (\mlZeroFspCi; healthy valid actions \mlZeroHealthy; heads selected on this same cohort, a screen not yet reviewed), and adding the full task vector of TMax-4B, an RL-trained Qwen3.5-4B and the only one of four planned RL arms that was run, by \rlFsp{} (\rlFspCi; healthy valid actions \rlHealthy, failing the healthy gate). % source: data/pod_astra/moonleak/screen_head4_s0.json; data/pod_astra/moonrl/analysis_rl100.json, REVIEW_RL100.md (figs/power.py)
All intervals are task bootstraps with \pwBootReps{} resamples, uncorrected for the four contrasts.

\subsection{Online evaluation}
\label{sec:online}
No candidate passed the gate, so none was compiled into the weights (Appendix~\ref{sec:compile} states the compilation step, Proposition~\ref{prop:equiv}); five arms were run online, each paired by task with an original run on the same harness, and the arms served by Modal with an original also served by Modal, on the same GPU type, image and vLLM version but from a separate account (Table~\ref{tab:b4-heldout}). The \texttt{answer} output-row edit changes only the row of the (untied) output head for the first token of \texttt{answer} (2{,}560 weights), lowering its logit by about 40 at the tool-name positions where it was calibrated; the corrective prompt adds one line to the end of the system prompt: ``Before calling answer, look at the current screenshot and check that every requirement of the task is visibly satisfied; if anything is missing or unconfirmed, continue working instead of answering.'' A sixth arm, the output head scaled by 1.6 (an effective sampling temperature of 0.5), was started on the reserve and the \textsc{test} split, but its server froze from 00:44 to 03:22 UTC, so both runs were invalidated and are not reported. % source: ../JOURNAL.md 21:44 UTC (eager_lm40: row 8944 of lm_head, c=0.26, answer logit -40.0 at calibration positions); osworld/run_holo_verify.py VERIFY (prompt v2_verify in meta.json of results_val_prompt_eager_m40); ../JOURNAL.md 03:25 and 03:42 UTC (lmscale16 INVALID); scripts/make_patched_ckpt.py lm_head_scale The original itself is the first measure of noise: its three runs on the confirmation reserve, spread over \onReserveOrigRange{} points (first row of Table~\ref{tab:b4-heldout}), are the run-to-run noise against which every arm is read. It scores \onReserveOrigPct\% on one harness and, on the other, \onReserveGhOrigPct\% when served by our H200 server and \onReserveModalOrigPct\% when served by Modal: the same weights differ by \onNoiseDelta{} points across harnesses (\onNoiseDisc{} tasks better and worse) and by \onNoiseServerDelta{} points across servers (95\% CI \onNoiseServerCi), differences comparable to the effects below, every one of which lies within a 95\% CI that includes zero. Against it, no arm shows a gain: ExPO at $\lambda=1.25$ \onReserveExpoDelta{} points (95\% CI \onReserveExpoCi), the \texttt{answer} output-row edit \onReserveLmrowArm{} against \onReserveLmrowOrig, the corrective prompt \onReservePromptArm{} against \onReservePromptOrig, and, against the Modal-served original, the CUA-Lite merge \onReserveCuaDelta{} (\onReserveCuaCi; \onReserveCuaDisc{} tasks better and worse) and the closed-form edit \onReserveBroadDelta{} (\onReserveBroadCi). % source: data/local/online/ (figs/online.py)
Two reserve tasks score 0 under the H200-served original and 1 under both Modal-served arms; they score 0 under the Modal-served original as well, so the gap came from run-to-run variation, not from the server. The harnesses do differ: two \texttt{os} tasks score 0 on the local harness and 1 on the GitHub one, so runs from different harnesses are never pooled. The merge also produced 6 unparsable steps against none for the H200-served original, a sign that it degrades the output format. % source: data/local/online/ (figs/online.py, SERVER_SENSITIVE asserted 0 under results_confirm_gh_orig_modal_confirm); data/local/reports/REDTEAM2_0335.md point 2
The decision rule for the merge was fixed before the comparison with the Modal-served original was read, but after the merge's own reserve scores had been seen, so only its \textsc{test} part is strictly prospective: the merge and an original served by Modal are compared on the 45 reserve and 120 \textsc{test} tasks, and a gain is claimed only if the lower bound of the task-bootstrap 95\% CI of the paired difference is above zero. By this rule, the merge shows no detectable gain: over the \onAllCuaN{} paired tasks it scores \onAllCuaArm{} against \onAllCuaOrig{} ($\Delta$ \onAllCuaDelta{} points, 95\% CI \onAllCuaCi; \onAllCuaDisc{} tasks better and worse), and on the \textsc{test} split alone, the strictly prospective part, \onTestCuaArm{} against \onTestCuaOrig{} ($\Delta$ \onTestCuaDelta, CI \onTestCuaCi). On the same \textsc{test} tasks the H200-served original scores \onTestOrigScore/120, \onTestNoiseServerDelta{} points above the Modal-served one (CI \onTestNoiseServerCi): the noise between two servings of the same weights is as large as the merge's effect. % source: data/local/online/results_test_gh_merge_cua_b050, results_test_gh_orig_modal_test, results_confirm_gh_*_modal_confirm (figs/online.py)
ExPO was also run on the \textsc{test} split without having passed the gate, a decision of the LLM agent coordinating the study, reported as such: \onTestExpoArm{} against \onTestExpoOrig{} ($\Delta$ \onTestExpoDelta{} points, 95\% CI \onTestExpoCi; \onTestExpoDisc{} tasks better and worse). All online scores are the official OSWorld score, a mean over tasks with partial credit included. With few tasks scoring differently between two runs (for the merge on the reserve, \onReserveCuaDisc{} better and worse), the bootstrap bounds move by up to a few points with the resampling seed; we report seed 0 throughout, and an LLM-agent recomputation agrees on every point estimate. Counting only full successes, the \textsc{test} difference for ExPO is \onTestExpoBinDelta{} points (95\% CI \onTestExpoBinCi). % source: data/local/online/test_compare.json (online-eval)
A paired difference smaller than about \onMde{} points would go undetected with 80\% power for ExPO on the reserve; the same computation gives about 21 and 23 points for the merge and the closed-form edit on the reserve, 50 for the output-row edit, 10 for each \textsc{test} comparison and 9 for the 165-task primary comparison, so these comparisons exclude only large effects. % source: (z_.975+z_.80)*SD(paired diff)/sqrt(n) as in figs/online.py onMde, on data/local/online runs: reserve CUA 20.7, closed-form 22.7, lm40 49.6 (n=14), TEST ExPO 10.5, TEST CUA 9.8, all CUA 9.1 (n=165)
A cheaper paired design, branching from recorded episodes, was piloted: replaying a recorded episode up to its last step reproduces the screenshots in \brValidOk{} of \brValidN{} validation episodes, and from the final step of \brEpisodes{} recorded episodes (two continuations each, both arms served by Modal) the original succeeds in \brOrig{} continuations and the merge in \brCand, with one episode whose replay diverges from its first step. The only failed episode that an original continuation recovers, \texttt{\brRepairDomain/\brRepairTask}, is an officially infeasible task, where the continuation states that the requested theme does not exist; on the feasible false successes whose replay is faithful, no continuation succeeds (original \brFeasFalseOrig, merge \brFeasFalseCand). The pilot shows no advantage for the merge, no repair of a feasible false success, and is too small to show anything else. % source: data/local/online/branch_pilot/REPORT.md and {orig,cand}/branches.jsonl (figs/online.py)
\placefig{tab-b4-heldout}
